\section{Introduction}

Transnational education (TNE) continues to expand across the UK higher education sector, with partnerships, validation, franchise delivery, and distance-learning arrangements forming a significant component of the international footprint of UK degree-awarding bodies. As provision grows in scale and complexity, the need for clear alignment between qualification levels, award structures, regulatory expectations, and partner-delivery responsibilities becomes increasingly important. Ensuring that programmes delivered in Sri Lanka under UK awarding arrangements remain compliant, transparent, and intelligible to students, regulators, and employers requires an evidence-based synthesis of frameworks, standards, and oversight models across the UK home nations.

The UK qualifications architecture is defined primarily by the Frameworks for Higher Education Qualifications (FHEQ) in England, Wales and Northern Ireland, and by the Scottish Credit and Qualifications Framework (SCQF) in Scotland. These frameworks set out level descriptors, typical credit volumes, award types, learning outcomes, and progression expectations. In Sri Lanka, the Sri Lanka Qualifications Framework (SLQF) provides the national reference point for qualification levels and credit norms, informing local recognition and programme evaluation.

At the regulatory level, England operates a registration-based system under the Office for Students (OfS), with conditions of registration that govern access and participation, quality and standards, financial sustainability, consumer protection, transparency, and information duties. Scotland, Wales, and Northern Ireland maintain distinct but related oversight models through the Scottish Funding Council (SFC), the Commission for Tertiary Education and Research (CTER), and the Department for the Economy (DfE~NI), respectively. UK~ENIC (ECCTIS) provides authoritative recognition advice for admissions, comparability, and credit transfer.

For UK providers delivering through Sri Lankan partners, it is essential that programme design aligns with FHEQ level descriptors, sector-recognised standards, and the expectations of UK regulators. This includes proportionate oversight of teaching, assessment, student outcomes, resources, support, student protection arrangements, and reportable events handling. Exit awards, progression routes, and credit volumes must be transparent and consistent to support recognition both locally and internationally.

This review synthesises the evidence base across levels, awards, governance, recognition, and TNE oversight to provide a coherent foundation for designing and assuring UK--Sri Lanka TNE provision. The outputs support the implementation of academically credible, regulatorily compliant, and student-centred programmes that remain intelligible to students, regulators, and employers.
